\section{Moment estimates for local random lists}\label{sec:norms}
The exact map of Section~\ref{sec:blocking} is defined for every
realization of the disorder.  We now estimate that map in probability.
There are two different tasks.  Low moments must agree, to a prescribed
accuracy, with the Taylor calculation at zero background.  A high moment
must remain small even though individual principal fields are unbounded.
The first task uses the rarity of a large field.  The second also uses its
geometric isolation.  Keeping these tasks separate will allow the second
moment tensors to control sup norms without a circular argument.

Write $X_i=\norm{h_i}_\infty$, with $n_i$ the slot length, and use the
weighted norms of Section~\ref{sec:blocking}.  All randomness below includes
the independent seeds reserved for local decisions.  In particular,
functions with disjoint enlarged dependence carriers are independent;
no assertion about independence of spins is being made.  Fix a finite set
of even source powers $q$, containing $2,4,400$, and an even integer
$p>8\max q$.  Increasing $p$ later in the initial choice is harmless.
Weights are chosen with the slack prescribed in Section~\ref{sec:blocking}:
even their fractions by $p$ and by the source powers exceed the animal
counting exponent.  Only finitely many weights occur.  They are fixed
before $L$, and compression of input lengths restores each of them in a
single step.

Here are the background hypotheses, for $0<s<1$:
\begin{equation}\label{norm:input}
 \norm{h}_{v,p}\le Bs,\qquad \norm{h}_{v,2}\le B_2s,
 \qquad B\ge B_2\ge1.
\end{equation}
The deterministic tail bound remains part of admissibility.  Constants in
$O(\cdot)$ may depend on $B,B_2$, on a mixed-moment constant introduced
below, and on all the fixed geometric choices.  The notation $o(s)$ will
always mean $O(s^{1+\kappa})$ for some fixed $\kappa>0$, not a limit
uniform in these constants.

\subsection{The probability and weight estimates used in Taylor expansions}
For a principal slot $c$, let $\mathcal B_c$ be its bad event.  Since its
random threshold lies in $[\delta,2\delta]$,
\begin{equation}\label{norm:badprob}
 \Prob(\mathcal B_c)\le \delta^{-p}\E X_c^p\le C(Bs/\delta)^p.
\end{equation}
We first record the precise way in which this supplies additional powers
of $s$.  If $a_1,\ldots,a_k$ are nonnegative integers with
$D=\sum a_\ell\le p$, then H\"older's inequality gives
\begin{equation}\label{norm:product}
 \E\prod_{\ell=1}^k X_{i_\ell}^{a_\ell}
 \le (Bs)^D\exp\{-v\sum_\ell a_\ell n_{i_\ell}\}.
\end{equation}
If a bad indicator at $c$ is also present, use
$\one_{\mathcal B_c}\le(X_c/\delta)^{p-D}$ before applying H\"older.
For $D\le4$ this gives
\begin{equation}\label{norm:flagproduct}
 \E\left[\prod_\ell X_{i_\ell}^{a_\ell}
                  \one_{\mathcal B_c}\right]
 \le C\delta^{-(p-D)}(Bs)^p
       \exp\{-v\sum_\ell a_\ell n_{i_\ell}\}.
\end{equation}
These inequalities allow repeated indices and require no independence.

We will use Lemma~\ref{block:tree} in the following explicit manner.
At a vertex with $b$ children, reserve a length weight $e^{-a n}$ and use
\[
 (1+n)^b e^{-a n}\le C_a^{b+1}b!.
\]
The remaining length weight pays the animal choices and the requested
output weight, by Lemma~\ref{block:carrier}.  Ordering the children
absorbs $b!$; the number of rooted ordered trees of order $k$ is at most
$4^{k-1}$.  Consequently a coefficient with $k$ occurrences has summed
majorant $C_L^k$ times its numerical vertex bounds.  Distinguishing at
most $p+4\max q$ vertices multiplies this by a fixed power of $k$, which
can be absorbed by increasing $C_L$.  A bad test can be chosen in at most
$C_L\sum(1+n_i)$ positions, since every decision is local.  The same
reserved length weight pays this factor.  For tensor products there are
at most four output roots here; one applies this calculation to each
root before multiplying.  Thus all subsequent uses of H\"older are made
\emph{before} sums over indices.  In particular the exponential decrease
in every retained index is preserved.

\begin{proposition}[Taylor transport of background moments]\label{norm:taylor}
Under \eqref{norm:input}, the mean and the raw moment and cumulant arrays
with at most four arguments are given by the ordinary zero-background
Taylor rule, retaining total background degree at most four, with error
$O(s^5)$ in the prescribed weighted indexed projective norms.  The
coefficient maps are bounded uniformly in the scale.  An output cumulant
vanishes if its dependence carriers split into two nonempty mutually
separated families.
\end{proposition}
\begin{proof}
Expand each output slot into its pure bad part, its singleton part, and
its absolutely convergent connected small-field series.  The pure bad
part has centered sup norm bounded by twice the sum of the norms of the
bad principal inputs that generate it.  Each such input carries its own
bad indicator.  All conditional moments in the other two parts are
bounded by the product of the input sup norms, because the conditional
laws, including the tilted laws, are normalized.

First consider terms without a bad part.  On the event that all local
bad tests relevant to a fixed coefficient fail, that coefficient is the
ordinary Taylor coefficient.  The difference between the two versions
is bounded by their sum of absolute majorants and a union of their local
bad events.  For total degree $D\le4$, \eqref{norm:flagproduct} bounds
each such difference by $O(s^p)$ with all retained index weights.
The tree calculation above sums the coefficients and the possible test
centers.  A term containing a pure bad part is treated by expanding its
finite sum of principal input norms and choosing one of its bad
indicators.  For products of at most four outputs the same inequality
applies; each remaining factor either contributes an input norm or a
convergent small-field series.

It remains to bound degrees at least five in those series.  Retain five
input factors, use \eqref{norm:product} on them, and bound every other
good occurrence by its deterministic threshold or tail bound.  The
majorant is bounded by
\[
 C s^5\sum_{k\ge5} k^C C_L^k\delta^{k-5}=O(s^5),
\]
after decreasing $\delta$.  The same argument bounds a high-degree term
that has already supplied a bad part: retain enough input factors to
reach degree five, or use the bad indicator when fewer factors are
present.  There are at most four unrestricted bad-only factors, so
$p>8\max q$ supplies all required H\"older powers.  This proves the raw
moment assertion.

For clarity, the projective estimate here is an estimate before any
spin evaluation: a random elementary tensor has projective norm equal
to the product of its factor norms, and the norm of its expectation is
at most the expectation of that product.  Applying the finite partition
formula
\[
 \cum(Z_1,\ldots,Z_k)
 =\sum_{\pi\in\mathcal P([k])}(-1)^{|\pi|-1}(|\pi|-1)!
       \bigotimes_{A\in\pi}\E\bigotimes_{a\in A}Z_a
\]
with factors put in their original argument order proves the cumulant
assertion and its error bound.  Finally, disjoint dependence carriers
use disjoint primitive disorder variables and seeds.  Their joint moment
generating polynomial factors, so its logarithm has no coefficient
involving both families.  This is the claimed connected-index restriction.
\end{proof}

One related consequence, needed below, keeps the probability norm rather
than taking expectation.  Retaining two field factors and using
H\"older in $L^2$ gives an $O(s^2)$ bound for nonlinear good terms.
For a modified singleton or a bad-only term use
\[
 \norm{X_i\one_{\mathcal B_c}}_2
 \le \delta^{-(p-2)/2}
      \big(\E X_i^2X_c^{p-2}\big)^{1/2}
 \le C s^{p/2}e^{-v n_i}.
\]
Together with the same tree sums, this proves
\begin{equation}\label{norm:quadratic-L2}
                  \norm{h'-Th}_{v,2}=O(s^2).
\end{equation}
This estimate concerns the error in replacing the map by its linear
part.  It does not yet bound the $L^2$ sup norm of that linear part by
the second moment tensor; that is the role of the finite dictionary.

\subsection{Attenuation of an isolated positive tilt}
The estimate that handles unbounded principal fields is deterministic.
It uses the reference pointwise comparison before any disorder average.

\begin{lemma}[Positive-tilt attenuation]\label{norm:attenuation}
Let $B$ and $F$ be even functions supported in a square of radius $w$,
with a proportional empty collar, and let $P_R$ average to a surrounding
square of radius $R\ge Cw$.  There is a constant $C$, independent of the
real function $B$, such that
\begin{align}
 \inf_{a\in\R}\norm{\frac{P_R(e^BF)}{P_R e^B}-a}_\infty
       &\le C\frac wR\norm{F}_\infty,\label{norm:tilted-function}\\
 \inf_{a\in\R}\norm{\log P_R e^B-a}_\infty
       &\le C\frac wR\min\{1,\norm{B}_\infty\}.
                                      \label{norm:tilted-log}
\end{align}
The centered versions satisfy the same bounds after multiplying $C$ by
at most two.  In particular, $B$ need not define a ferromagnetic
perturbation or satisfy a smallness assumption.
\end{lemma}
\begin{proof}
Choose an intermediate square, of radius a fixed multiple of $w$, whose
contour lies beyond the collar.  Put
\[
 D=\frac{P_{Cw}e^B}{\langle e^B\rangle_0},\qquad
 N=\frac{P_{Cw}(e^BF)}{\langle e^B\rangle_0}.
\]
Lemma~\ref{ref:comparison} bounds, pointwise, the reference law of
the inner spins under any contour assignment above and below by fixed
multiples of its plane law.  Integration against $e^B>0$ therefore gives
$c\le D\le C$ and $|N|\le C\norm F_\infty$, uniformly in $B$.
Both functions are even and $\langle D\rangle_0=1$.
The even part of Proposition~\ref{ref:operator}, followed by the tower
property, gives
\[
 P_RD=1+d_R,\qquad P_RN=\langle N\rangle_0+n_R,
 \qquad \norm{d_R}_\infty\le Cw/R,
 \quad \norm{n_R}_\infty\le C(w/R)\norm F_\infty.
\]
Since $P_RD\ge c$, subtracting $\langle N\rangle_0$ from their ratio
proves \eqref{norm:tilted-function}.  The estimate remains uniform upon
replacing $B$ by $tB$, $0\le t\le1$.  Integrate
\[
 \frac{\dd}{\dd t}\log P_Re^{tB}
       =\frac{P_R(Be^{tB})}{P_Re^{tB}}
\]
in the quotient by constants and apply \eqref{norm:tilted-function}
with $F=B$.  This proves \eqref{norm:tilted-log} without introducing
$e^{\norm B}$ for the bound linear in $\norm B_\infty$.
For the bound independent of that norm, subtract
$\log\langle e^B\rangle_0$ and use
$\log P_RD=\log(1+d_R)$: since $P_RD\ge c$, its norm is at most
$C\norm{d_R}_\infty\le Cw/R$.  Taking the smaller of the two bounds
proves the stated minimum.
\end{proof}

Set $M=(\log L)^3$ and enlarge $L$ until $M$ dominates all fixed halos.
For every short input, meaning $n_i\le M$, the geometric bad-component
rule has the following consequence.  Its contribution to the absolute
bound for a principal output may be charged
\begin{equation}\label{norm:witness-bound}
       C\frac ML X_i
       +C X_i\sum_{c\in\mathcal W(i)}\one_{\mathcal B_c}.
\end{equation}
Here $\mathcal W(i)$ is a deterministic set of at most $C_L$ principal
centers within $CL$ fine cells of the input anchor and at distance more
than $C_0M$ from it.  The fixed $C_0$ may be chosen large enough that a
witness carrier is disjoint from this input and from any second short
input whose anchor is within $10M$ of it.

We verify the geometric implication, including the case of a large
input.  If the singleton is unaffected by bad centers, apply the even
reference estimate on radius $CMr$ and then the tower property.  If it
is affected and no witness occurs, all affecting bad centers lie in a
square of radius $CMr$ about its anchor.  Their contact component is
selected: otherwise the first contact leaving this square supplies a
bad center in the decision window, and hence a witness.  The component,
the singleton support, and their halos consequently lie deep inside the
selected cutting square.  Lemma~\ref{norm:attenuation} applies to their
union.  For a good singleton use \eqref{norm:tilted-function}; for the
sum of the bad inputs use \eqref{norm:tilted-log} and then
$\norm{\sum b}_\infty\le\sum\norm b_\infty$.  If a selected square is
nearby but its bad support is farther away than the inner square just
described, one of those bad centers is itself a witness.  On the witness
event use the contraction of a normalized conditional expectation, or
the bound $2\sum\norm b_\infty$ for a centered logarithm.  This proves
\eqref{norm:witness-bound}.  Enlarging the numerical distances in the
blocking convention makes $C_0$ compatible with the component cutoff;
all constants remain independent of $L$.
Figure~\ref{norm:geometry} shows the inner collar used for attenuation
and the separated dependence carriers used on the witness event.
\begin{figure}[tbp]
\centering
\begin{tikzpicture}[x=1cm,y=1cm,font=\small,
  outer/.style={draw=blue!55!black,line width=.8pt},
  inner/.style={draw=blue!55!black,line width=.65pt},
  footprint/.style={draw=black!65,densely dashed,line width=.65pt},
  guide/.style={draw=black!45,densely dotted},
  dimension/.style={draw=black!65,{Stealth[length=1.5mm]}-{Stealth[length=1.5mm]}}]
  \node[anchor=south,font=\small\bfseries] at (0,2.75)
    {(a) Isolated bad component};
  \draw[outer] (-2.5,-2.5) rectangle (2.5,2.5);
  \fill[blue!5] (-1.65,-1.65) rectangle (1.65,1.65);
  \fill[white] (-1.05,-.9) rectangle (1.05,1.2);
  \draw[inner] (-1.65,-1.65) rectangle (1.65,1.65);
  \draw[guide] (-1.05,-.9) rectangle (1.05,1.2);
  \node[align=center] at (0,.35) {Bad supports\\and input $h_i$};
  \node[font=\footnotesize] at (0,-.57) {radius $O(Mr)$};
  \node[fill=blue!5,inner sep=1.5pt] at (0,-1.27) {Proportional collar};
  \node[fill=white,inner sep=1.5pt] at (0,2.09)
    {Outer cut: radius $HR$};
  \node[fill=white,inner sep=1.5pt] at (0,-2.1)
    {Inner normalization square};
  \draw[black!65,-{Stealth[length=1.3mm]}] (0,-1.89) -- (0,-1.67);
  \node at (0,-2.90) {$R=Lr,\qquad M\ll L$};

  \begin{scope}[xshift=7.25cm]
    \node[anchor=south,font=\small\bfseries] at (0,2.75)
      {(b) A distant bad witness};
    \draw[guide] (-2.5,-2.5) rectangle (2.5,2.5);
    \node[fill=white,inner sep=1.5pt] at (0,2.08)
      {Local decision window: $O(Lr)$};
    \draw[footprint,fill=blue!4] (-2.0,-.65) rectangle (-.7,.65);
    \fill (-1.35,0) circle[radius=1.5pt];
    \node[anchor=south] at (-1.35,.1) {$c_i$};
    \node[anchor=north,align=center] at (-1.35,-.77)
      {Input footprint\\$n_i\le M$};
    \draw[footprint,fill=orange!7] (.90,-.42) rectangle (1.74,.42);
    \fill[orange!65!black] (1.32,0) circle[radius=1.8pt];
    \node[anchor=west,inner sep=1pt] at (1.43,0) {$z$};
    \node[anchor=north,align=center] at (1.32,-.77)
      {Flag footprint};
    \draw[dimension] (-1.35,1.1) -- (1.32,1.1)
      node[midway,above,inner sep=2pt] {distance $>C_0Mr$};
    \draw[guide] (-1.35,.65) -- (-1.35,1.25);
    \draw[guide] (1.32,.42) -- (1.32,1.25);
    \node[align=center] at (0,-2.04) {Independent randomness};
    \node[font=\footnotesize] at (0,-2.90) {Schematic; not to scale.};
  \end{scope}
\end{tikzpicture}
\caption{The two geometric mechanisms in the short-input estimate.
In (a), the supports of the bad fields and the short input lie in the
dotted square of radius $O(Mr)$.  The middle solid square provides a
proportional collar for pointwise
normalization of the positive bad weight before averaging to the
selected outer cut of radius $HR$.  The ratio of these radii produces
the $O(M/L)$ centered attenuation.  In (b), an exceptional short-input
contribution has a bad-center witness $z$ farther than $C_0Mr$ from
its anchor $c_i$, but within the local decision window of size
$O(Lr)$, where $C_0$ is the fixed witness-separation constant.
The separated footprints include the disorder and all
decision seeds, so the witness flag is independent of that input.
The panels illustrate separate cases; no particular arrangement of
bad fields or grid cells is prescribed.}
\label{norm:geometry}
\end{figure}


Long singletons, $n_i>M$, have a separate small factor.  Reserve an
input length weight after paying the output weight.  Their summed
majorant is at most
\begin{equation}\label{norm:long}
           \lambda_L=C L^C e^{-cM},
\end{equation}
times the relevant input norm.  The constants in the exponents depend
only on the already fixed weights and powers.  In particular $\lambda_L$
is smaller than every inverse power of $L$ as $L\to\infty$.
Finally, an analytic background term of order at least two is bounded
by $C_L\delta$ times a positive linear sum of input $X_i$: retain one
factor, bound all others by their good-field bounds, and sum the trees.
This observation holds in any probability measure, including a measure
weighted by a source.

\subsection{Reproduction of the background probability norms}
\begin{proposition}[High moments and recovery of $L^2$]\label{norm:reproduction}
Choose $L$ sufficiently large and then $\delta$ sufficiently small.
Under \eqref{norm:input},
\begin{align}
 \norm{h'}_{v,p}&\le Bs/8+C_L^*B_2s+o(s),\label{norm:high-reproduction}\\
 \norm{h'}_{v,2}&\le C_{\eta,L}F_*s+\eps_* Bs+o(s).
                                                    \label{norm:L2-reproduction}
\end{align}
Here $C_L^*$ grows at most polynomially in $L$.  For the second inequality,
assume that each output principal slot has diagonal second-moment
bound
\begin{equation}\label{norm:diagonal}
 \sup_\sigma\E|h'_a(\sigma)|^2\le e^{-2v}F_*^2s^2,
 \qquad F_*\ge1.
\end{equation}
A bound on its raw second tensor implies this hypothesis.  One can
choose $\eps_*<[100(1+C_L^*)]^{-1}$: first use the long-slot suppression,
then choose the dictionary accuracy $\eta$ and the mesh threshold $r_0$.
These choices do not depend on $B,B_2,F_*$.
\end{proposition}
\begin{proof}
For the first inequality, use the absolute decomposition just proved.
The nonlinear good terms cost $C_L\delta Bs$, and long singleton terms
cost $\lambda_LBs$.  A witness is independent of the short input it
multiplies.  Thus
\[
 \norm{X_i\one_{\mathcal B_c}}_p
  =\norm{X_i}_p\Prob(\mathcal B_c)^{1/p}
  \le C B^2\delta^{-1}s^2e^{-vn_i}.
\]
Minkowski's inequality and the finite witness sums give $O(s^2)$ at
fixed $L$.  It remains to estimate the attenuated short sum.

Expand its $p$th power.  Join two of the $p$ input anchors when their
distance is at most $C M$, choosing $C$ so that different components
have disjoint dependence carriers.  Call the connected components
bunches.  For a term with one bunch, choose the first anchor in at most
$CL^2$ ways.  A spanning tree on the $p$ labeled occurrences places
each successive anchor in at most $CM^2$ ways.  There are at most
$p^{p-2}$ such trees.  Animal shapes and slot lengths are summed with
the reserved weights.  Hence the count is at most $CL^2M^{C_p}$, and
\eqref{norm:product} bounds its expectation by $(Bs)^p$.  After taking
the $p$th root and including attenuation, this part costs
\begin{equation}\label{norm:one-bunch}
                 C L^{-1+2/p}M^{C_p} Bs.
\end{equation}

For two or more bunches, their random variables are independent.
Suppose a bunch contains $k<p$ occurrences.  H\"older within it uses
only $L^k$ norms.  For $2\le k<p$, interpolate by
\[
 \norm{X_i}_k\le\norm{X_i}_2^{\theta_k}
                   \norm{X_i}_p^{1-\theta_k},
 \qquad \theta_k=\frac{2(p-k)}{k(p-2)};
\]
for $k=1$ use $\norm{X_i}_1\le\norm{X_i}_2$.  The product of all
bunch estimates, after its $p$th root, is therefore
\[
 s B^{1-\theta}B_2^\theta
   \exp\{-v\textstyle\sum n_i/p\},\qquad \theta\ge1/p.
\]
For the last inequality, a singleton already contributes one to
$p\theta$.  If there are no singletons and there are $b\ge2$ bunches,
then $\sum k\theta_k=2p(b-1)/(p-2)\ge1$.  All placement
counts are polynomial in $L$: there are at most $p$ bunch roots, each
with $CL^2$ choices, followed by the preceding short-distance counts.
For each of the finitely many bunch patterns, weighted Young's
inequality gives, for any fixed $\varepsilon>0$,
\[
 C L^C B^{1-\theta}B_2^\theta
       \le\varepsilon B+C_{p,\varepsilon}L^{C_p}B_2.
\]
This also follows directly by maximizing the left side minus
$\varepsilon B$ as a function of $B/B_2$.  Since $1/\theta\le p$,
the resulting constant remains polynomial in $L$.  Sum over the
finitely many patterns.  Increase $L$ to make \eqref{norm:one-bunch}
and \eqref{norm:long} small; choose $\varepsilon$ for the multi-bunch
terms; then decrease $\delta$ for the nonlinear terms.  Allocating
less than $1/8$ in total to these coefficients of $B$ proves
\eqref{norm:high-reproduction}.

For the second inequality, first use \eqref{norm:quadratic-L2}.
The nonprincipal part of $Th$ has input length at least $L/2$, and its
$L^2$ norm is at most $CL^Ce^{-cL} Bs$.  A principal output of $Th$
has the form $\pi P_{HR}A$, where $A$ is the sum of the short-rule
inputs anchored in its parent cell and $\pi$ denotes plane centering.
All these inputs are supported in one inner square, separated by a
fixed proportional buffer from the $HR$ contour.  The weighted
placement bound and Minkowski give
\[
                     \norm{\norm A_\infty}_2\le C_L Bs.
\]
The finite dictionary of Lemma~\ref{ref:dictionary} supplies $N=N(\eta,L)$
actual contour assignments $\sigma_1,\ldots,\sigma_N$ such that
\[
 \norm{\pi P_{HR}A}_\infty
 \le \max_{\nu\le N}|\pi P_{HR}A(\sigma_\nu)|
                          +C\eta\norm A_\infty.
\]
This is a deterministic inequality and thus applies to the random $A$.
Replace each tested value of $Th$ by that of $h'$ using
\eqref{norm:quadratic-L2}.  The $L^2$ norm of the maximum of $N$ values
is at most the square root of the sum of their second moments.  By
\eqref{norm:diagonal}, after restoring the principal weight, this costs
$\sqrt N F_*s+O(s^2)$.  The net error costs $C_L\eta Bs$.
Choose $L$ so that the exponential long-slot coefficient is smaller
than the desired $\eps_*/2$; this is possible against the polynomial
$C_L^*$.  Then choose $\eta$ with $C_L\eta<\eps_*/2$ and increase
$r_0$ to the dictionary threshold.  This proves
\eqref{norm:L2-reproduction}.  Notice the order of the argument:
\eqref{norm:diagonal} is a tensor hypothesis, supplied later by the
Taylor flow, whereas the conclusion is a probability norm of a sup
norm.  They are not the same assertion.
\end{proof}

\subsection{Sources and mixed probability measures}
Let $u$ be a scaled singleton odd source list.  Put $Y_i=\norm{u_i}_\infty$
and suppose, for one of the chosen even $q$, that
\begin{align}
 (\E Y_i^q)^{1/q}&\le Ue^{-vn_i},\label{norm:source-input}\\
 (\E Y_i^qX_m^p)^{1/p}
  &\le U^{q/p}B_s s\,e^{-vq n_i/p-vn_m}.
                                      \label{norm:mixed-input}
\end{align}
The subscript on $B_s$ denotes ``source'': $B_s$ is a fixed constant,
whereas $U$ may change from one scale to the next.  No independence
between $u_i$ and a nearby $h_m$ is assumed.

\begin{proposition}[Source Taylor estimates and mixed reproduction]
\label{norm:source}
Under the preceding assumptions, moments with $q$ source occurrences
and at most two background arguments obey the ordinary Taylor rule
through total background degree two,
with error $O(U^qs^3)$ in the weighted projective norms.  Background
arguments already present in a tensor count toward that degree.
With fewer source arguments the corresponding bound follows by
H\"older.  Moreover,
\begin{equation}\label{norm:source-difference}
               \norm{u'-L^dTu}_{v,q}=O(Us),
\end{equation}
and the output satisfies
\begin{equation}\label{norm:mixed-output}
 \sup_{a,b}e^{vq n_a/p+vn_b}
       (\E\norm{u'_a}_\infty^q\norm{h'_b}_\infty^p)^{1/p}
 \le U^{q/p}\{B_s s/8+C_{L,B}s+o(s)\}.
\end{equation}
The choices of $p,L,\delta$ and the weights can be made before
$B_s$ and $U$.  All these estimates retain any saved scalar's carrier
until after its weighted estimate is taken.
\end{proposition}
\begin{proof}
For a fixed source input define the finite measure
$\dd\mu_i=Y_i^q\dd\Prob$.  Its mass is at most
$U^qe^{-vq n_i}$, and \eqref{norm:mixed-input} bounds the $p$th moment
of every $X_m$ in this measure.  Consequently, for nonnegative integers
$a_m$ of total $D\le p$,
\begin{equation}\label{norm:weighted-holder}
 \int\prod_m X_m^{a_m}\,\dd\mu_i
 \le U^q(B_s s)^D
       e^{-vq n_i-v\sum_m a_mn_m}.
\end{equation}
Indeed, apply H\"older with exponents $p/a_m$ and $p/(p-D)$, the last
factor being the mass of $\mu_i$.  The powers of $U$ and of the source
weight sum to exactly $q$.  A bad indicator supplies the unused
$p-D$ powers as in \eqref{norm:flagproduct}.

For $q$ possibly distinct source slots, H\"older first gives
\[
 \E\left[\prod_{\ell=1}^qY_{i_\ell}\prod_mX_m^{a_m}\right]
 \le\prod_{\ell=1}^q
       \left(\int\prod_mX_m^{a_m}\,\dd\mu_{i_\ell}\right)^{1/q}.
\]
Thus \eqref{norm:weighted-holder} retains decay in every source and
background index.  For fewer source occurrences include unweighted
probability measures in this H\"older step.  Apply this calculation to
the tree expansion, retaining three background factors in a remainder.
Additional good factors are bounded deterministically, and a changed
bad decision is paid by the unused mixed moment powers.  The same
summable series as in Proposition~\ref{norm:taylor} proves the asserted
$O(U^qs^3)$ bound.  In $L^q$ retain one background factor and use
\eqref{norm:weighted-holder} with degree $q$.  A bad flag costs
$O(Us^{p/q})$.  This proves \eqref{norm:source-difference}, since
$p/q>1$.

We give the high mixed estimate separately.  Normalized tilted
conditional moments and the deterministic good-field bounds imply a
positive linear majorant
\begin{equation}\label{norm:source-majorant}
             Y'_a\le\sum_i K_{ai}Y_i,\qquad K_{ai}\ge0.
\end{equation}
Here the coefficients are deterministic majorants, not the realized
coefficients of the map.  They include the factor $L^d$.  Lemma
\ref{block:tree}, with the source vertex retained, proves that their
weighted row sums are bounded by $C L^d$; reserving additional input
length weight makes the part with $n_i>M$ bounded by
$CL^Ce^{-cM}$.  These statements mean row sums after paying the
requested output weight from the input weight; they also hold with
all the fractional weights appearing below.  In particular no factor
$\norm{\log P e^B}$ occurs in \eqref{norm:source-majorant}.
To construct $K_{ai}$, sum the absolute tree majorants over the allowed
carrier and anchor assignments, taking the prescribed enlarged union
of input carriers for each output.  Every realized routing is one of
these boundedly many assignments.  There is no sum over patterns of
threshold outcomes or over the amplitudes of bad fields.

Fix output background slot $b$ and source input $i$.  Estimate
$\norm{\norm{h'_b}_\infty}_{L^p(\mu_i)}$ using the deterministic
decomposition preceding Proposition~\ref{norm:reproduction}.  If
$n_i>M$, crude Minkowski and \eqref{norm:mixed-input} give
\begin{equation}\label{norm:mixed-long-source}
 \norm{\norm{h'_b}_\infty}_{L^p(\mu_i)}
 \le C_L U^{q/p}B_s s\,e^{-vq n_i/p-vn_b}.
\end{equation}
There is still a reserved exponential in $n_i$ for its subsequent
source sum.

Suppose now $n_i\le M$.  Long background inputs and nonlinear good
terms cost, respectively, $\lambda_L B_s s$ and $C_L\delta B_s s$,
times the common factor $U^{q/p}e^{-vq n_i/p-vn_b}$.
For short background inputs with anchors within $10M$ of the source
anchor, use \eqref{norm:mixed-input}.  There are at most $CM^2$
anchors.  Animal lengths and shapes cost a constant after their weight
sums.  The attenuated part of \eqref{norm:witness-bound}, estimated
by Minkowski, therefore costs
\[
       C M^C L^{-1} B_s s\,
                  U^{q/p}e^{-vq n_i/p-vn_b}.
\]
For its witness part the witness carrier is disjoint from the union
of the source and background input carriers.  Exact independence gives
\[
 \big(\E Y_i^q X_m^p\one_{\mathcal B_c}\big)^{1/p}
   =\big(\E Y_i^qX_m^p\big)^{1/p}\Prob(\mathcal B_c)^{1/p},
\]
so this part gains an additional factor $O(Bs/\delta)$ and is
$O(U^{q/p}s^2)$ at fixed constants.  For short background anchors
farther than $10M$, the background input is independent of $Y_i$.
Drop any witness indicator and use
\[
 \big(\E Y_i^qX_m^p\big)^{1/p}
       = (\E Y_i^q)^{1/p}(\E X_m^p)^{1/p}
       \le U^{q/p} Bs\,e^{-vq n_i/p-vn_m}.
\]
Even a crude sum over these anchors has only a polynomial cost in $L$.
Combining these estimates gives, for short $i$,
\begin{equation}\label{norm:mixed-fixed-source}
 \norm{\norm{h'_b}_\infty}_{L^p(\mu_i)}
 \le U^{q/p}e^{-vq n_i/p-vn_b}
       \{a_LB_s s+C_LBs+o(s)\},
 \quad a_L=C M^C/L+\lambda_L+C_L\delta.
\end{equation}

We now check the powers when summing sources.  Put
$F_i=(\E Y_i^q\norm{h'_b}_\infty^p)^{1/p}$.
Minkowski in $L^q$, applied to \eqref{norm:source-majorant}, yields
\begin{equation}\label{norm:mixed-minkowski}
 (\E (Y'_a)^q\norm{h'_b}_\infty^p)^{1/p}
             \le \left(\sum_i K_{ai}F_i^{p/q}\right)^{q/p}.
\end{equation}
This is an $\ell^{p/q}$ norm with weights $K_{ai}$.  Apply its triangle
inequality to the three terms in braces in
\eqref{norm:mixed-fixed-source}.  Raising the input source weight to
$p/q$ gives $(e^{-vq n_i/p})^{p/q}=e^{-vn_i}$, exactly the weight needed in
the positive source sum.  The background output weight remains
$e^{-vn_b}$ after taking the power $q/p$.  Hence the short-source
coefficient of $B_s s$ is at most
\[
          (CL^d)^{q/p}\big(CM^C/L+\lambda_L+C_L\delta\big),
\]
whereas the far-background term is a constant $C_{L,B}s$ independent
of $B_s$.  For the long-source part use
\eqref{norm:mixed-long-source} and the reserved exponential in its
row sum; its coefficient is at most $CL^Ce^{-cM q/p}$.
First choose $L$ large.  Since $dq/p<1$ and $M$ is a power of $\log L$,
the first attenuation coefficient tends to zero, as do both long
coefficients.  Then choose $\delta$ small.  Their sum can be made less
than $1/8$.  Witness errors have an extra positive power of $s$ and
remain $o(s)$ after \eqref{norm:mixed-minkowski}, whose use of the
$\ell^{p/q}$ triangle inequality is homogeneous of degree one in $F_i$.
This proves \eqref{norm:mixed-output}.
\end{proof}

The deterministic estimate \eqref{norm:source-majorant} also explains
why marked readout is possible in boxes whose principal background
fields are large near the boundary.  Before any probability estimate,
the homogeneous transport of a coefficient with a fixed set of mark
labels has a bounded weighted operator cost per step, with the
prescribed factors $L^d$ for its marks.  A freshly generated coefficient
with two distinct marks has a bounded bilinear cost in the two singleton
lists and retains both carriers and its saved scalar carrier; apply
Lemma~\ref{block:tree} with two source vertices.
A background factor retained at a merge can then be estimated under
the source-weighted measures above, even when the two sources are
correlated with it.  No large-field logarithm multiplies either source.

We finish by making the dependence of choices explicit.  Fix the
source powers, then $p$, the finite collection of weights and their
counting slack, and the geometric halos.  Choose $L$ for attenuation,
long-slot compression and the one-bunch bound.  Choose $\delta$ for
convergence and the small nonlinear coefficients.  The polynomial
$C_L^*$ is now fixed.  Choose $\eps_*$, dictionary accuracy $\eta$, and
$r_0$ in that order.  The moment constants $B_2,B,B_s$ and the tensor
allowances may then be fixed, and finally $s$ made sufficiently small.
In particular an allowance depending on the starting mesh does not
force a new mesh choice.  Because every contraction above has fixed
slack, the same inequalities hold on replacing the next-scale targets
$s,U$ by $s(1+O(s^2))$ and $U(1+O(s))$, respectively.  The stronger
variation $U(1+O(s^2))$ used after the marked cumulant comparison is
therefore covered as well.
